\documentclass[11pt,a4paper]{article}
\usepackage[utf8]{inputenc}
\usepackage{amsmath,amssymb,amsthm}
\usepackage{algorithm}
\usepackage{algpseudocode}
\usepackage{booktabs}
\usepackage{hyperref}
\usepackage{geometry}
\geometry{margin=1in}

\title{Truncated Backpropagation Through Time (TBPTT): Windowed Gradient Truncation, Streaming Memory, and Bias Bounds}
\author{Team Scaibu \\ \small Email: \texttt{jaydeep.9664920749@gmail.com} \\ \small Architecture and Core Engine Team}
\date{\today}

\newtheorem{theorem}{Theorem}
\newtheorem{lemma}[theorem]{Lemma}
\newtheorem{proposition}[theorem]{Proposition}
\theoremstyle{definition}
\newtheorem{definition}{Definition}
\newtheorem{remark}{Remark}

\begin{document}

\maketitle

\begin{abstract}
Full Backpropagation Through Time (BPTT) becomes computationally intractable and memory-prohibitive when training recurrent neural architectures on long continuous sequences or streaming time series ($T \to \infty$). Truncated BPTT (TBPTT) resolves this scaling bottleneck by partitioning sequences into localized chunks of window length $k_1$ forward steps and $k_2$ backward unrolling steps. By maintaining persistent hidden state vectors across chunks while severing the computational gradient graph at chunk boundaries, TBPTT achieves bounded $\mathcal{O}(k_2)$ memory footprint. This document formalizes the windowed gradient surrogate, analyzes truncation bias bounds, presents complete pseudocode, and walks through a verified numerical example.
\end{abstract}

\section{Notation and Mathematical Preliminaries}

Let an arbitrary sequence $\mathcal{X} = (\mathbf{x}_1, \dots, \mathbf{x}_T)$ be segmented into contiguous chunks of length $k$. For chunk $m \in \{0, \dots, \lfloor T/k \rfloor - 1\}$, the active time window spans $t \in [mk, (m+1)k - 1]$.

\begin{table}[h]
\centering
\caption{Mathematical Notation for Truncated BPTT}
\begin{tabular}{lll}
\toprule
\textbf{Symbol} & \textbf{Domain} & \textbf{Semantic Description} \\
\midrule
$k_1$ & $\mathbb{N}$ & Number of forward propagation steps between parameter updates \\
$k_2$ & $\mathbb{N}$ & Number of backward unrolling steps ($k_2 \ge k_1$, typically $k_1 = k_2 = k$) \\
$\mathbf{h}_{\text{init}}^{(m)}$ & $\mathbb{R}^{d_h}$ & Detached initial hidden state for chunk $m$: $\text{detach}(\mathbf{h}_{mk-1}^{(m-1)})$ \\
$\mathcal{L}^{(m)}$ & $\mathbb{R}$ & Segmented chunk loss: $\sum_{t=mk}^{(m+1)k - 1} \mathcal{L}_t$ \\
$\mathbf{g}_{\text{TBPTT}}$ & $\mathbb{R}^P$ & Truncated gradient estimator: $\sum_m \nabla_\theta \mathcal{L}^{(m)}$ \\
$\mathbf{g}_{\text{exact}}$ & $\mathbb{R}^P$ & True full-sequence BPTT gradient \\
\bottomrule
\end{tabular}
\end{table}

\section{Problem Definition and Core Concepts}

\subsection{3-Level Explanation}
\begin{enumerate}
    \item \textbf{Intuitive (High School level):} If you are reading an endless book, you can't remember every single letter from page 1 when you are on page 500. Instead, you keep your current train of thought, but only look back 20 words to fix recent grammar mistakes.
    \item \textbf{Engineering Level:} Full BPTT on $T=100{,}000$ steps requires storing $100{,}000$ hidden state activation vectors in GPU RAM before backward pass. TBPTT bounds memory to a fixed window $k$ (e.g., $k=64$), allowing training on arbitrarily long streams.
    \item \textbf{Mathematical Level:} TBPTT replaces the exact gradient $\nabla_\theta \mathcal{L} = \sum_{t=1}^T \sum_{\tau=1}^t \frac{\partial \mathcal{L}_t}{\partial \mathbf{h}_t} \frac{\partial \mathbf{h}_t}{\partial \mathbf{h}_\tau} \frac{\partial \mathbf{h}_\tau}{\partial \theta}$ with a truncated estimator:
    \begin{equation}
    \nabla_\theta \mathcal{L}_{\text{TBPTT}}(k_2) = \sum_{t=1}^T \sum_{\tau = \max(1, t - k_2 + 1)}^t \frac{\partial \mathcal{L}_t}{\partial \mathbf{h}_t} \left( \prod_{j=\tau+1}^t \frac{\partial \mathbf{h}_j}{\partial \mathbf{h}_{j-1}} \right) \frac{\partial \mathbf{h}_\tau}{\partial \theta}
    \end{equation}
\end{enumerate}

\subsection{5-Step Algorithmic Framework}
\begin{enumerate}
    \item \textbf{Chunk Partitioning:} Slice continuous input stream into blocks of length $k$.
    \item \textbf{Detached State Ingestion:} Initialize chunk $m$ with final hidden state $\mathbf{h}_{mk-1}$ from chunk $m-1$, detaching incoming autograd nodes.
    \item \textbf{Forward Unroll on Chunk:} Compute $\mathbf{h}_t, \mathbf{y}_t$ and per-step losses $\mathcal{L}_t$ for $t \in [mk, (m+1)k - 1]$.
    \item \textbf{Backward Adjoint Propagation:} Unroll backwards $k$ steps, halting backpropagation at the chunk boundary ($t = mk$).
    \item \textbf{Immediate Parameter Update:} Apply optimizer step $\theta \gets \theta - \eta \nabla_\theta \mathcal{L}^{(m)}$ and advance to chunk $m+1$.
\end{enumerate}

\section{Algorithm Specification}

\begin{algorithm}[H]
\caption{Truncated Backpropagation Through Time (TBPTT$(k, k)$)}
\label{alg:tbptt}
\begin{algorithmic}[1]
\Require Full sequence $\mathcal{X} = \{\mathbf{x}_t\}_{t=1}^T$, targets $\mathcal{Y} = \{\mathbf{y}_t^*\}_{t=1}^T$, chunk length $k$, parameters $\theta = (W_x, W_h, W_y, \mathbf{b}_h, \mathbf{b}_y)$.
\Ensure Updated parameters $\theta$.
\State $\mathbf{h}_{\text{carry}} \gets \mathbf{0}$
\For{$m = 0$ \textbf{to} $\lfloor T/k \rfloor - 1$}
    \State $\mathbf{h}_0 \gets \text{Detach}(\mathbf{h}_{\text{carry}})$
    \State $\mathcal{L}_{\text{chunk}} \gets 0$
    \For{$\tau = 1$ \textbf{to} $k$} \Comment{Forward chunk pass}
        \State $t \gets m \cdot k + \tau$
        \State $\mathbf{h}_\tau \gets \tanh(W_h \mathbf{h}_{\tau-1} + W_x \mathbf{x}_t + \mathbf{b}_h)$
        \State $\mathbf{y}_\tau \gets W_y \mathbf{h}_\tau + \mathbf{b}_y$
        \State $\mathcal{L}_{\text{chunk}} \gets \mathcal{L}_{\text{chunk}} + \mathcal{L}(\mathbf{y}_\tau, \mathbf{y}_t^*)$
    \EndFor
    \State $\mathbf{h}_{\text{carry}} \gets \mathbf{h}_k$ \Comment{Save final hidden state for next chunk}
    \State $\nabla_\theta \mathcal{L}_{\text{chunk}} \gets \text{BPTT\_Backward}(\{\mathbf{h}_\tau\}_{\tau=0}^k, \{\mathbf{y}_\tau\}_{\tau=1}^k, \dots)$
    \State $\theta \gets \theta - \eta \nabla_\theta \mathcal{L}_{\text{chunk}}$
\EndFor
\State \Return $\theta$
\end{algorithmic}
\end{algorithm}

\section{Theoretical Guarantees and Truncation Error Bound}

\begin{theorem}[Truncation Gradient Bias Bound]
\label{thm:tbptt_bias}
Let the recurrent transition matrix satisfy spectral norm $\|W_h\|_2 = \lambda < 1$. The approximation error between the exact full-sequence gradient $\mathbf{g}_{\text{exact}}$ and the truncated gradient $\mathbf{g}_{\text{TBPTT}}(k)$ decays exponentially with truncation depth $k$:
\begin{equation}
\|\mathbf{g}_{\text{exact}} - \mathbf{g}_{\text{TBPTT}}(k)\|_2 \le C \cdot \frac{\lambda^{k+1}}{1 - \lambda}
\end{equation}
where $C = \sup_t \left\| \frac{\partial \mathcal{L}_t}{\partial \mathbf{h}_t} \right\|_2 \sup_\tau \left\| \frac{\partial \mathbf{h}_\tau}{\partial \theta} \right\|_2$.
\end{theorem}

\begin{proof}
The exact gradient decomposes as $\mathbf{g}_{\text{exact}} = \mathbf{g}_{\text{TBPTT}}(k) + \mathbf{R}(k)$ where the residual tail is:
\begin{equation}
\mathbf{R}(k) = \sum_{t=k+1}^T \sum_{\tau=1}^{t-k} \frac{\partial \mathcal{L}_t}{\partial \mathbf{h}_t} \left( \prod_{j=\tau+1}^t \frac{\partial \mathbf{h}_j}{\partial \mathbf{h}_{j-1}} \right) \frac{\partial \mathbf{h}_\tau}{\partial \theta}
\end{equation}
Bounding the norm using $\|W_h\|_2 \le \lambda$:
\begin{equation}
\|\mathbf{R}(k)\|_2 \le C \sum_{t=k+1}^T \sum_{\tau=1}^{t-k} \lambda^{t-\tau} \le C \sum_{d=k+1}^\infty \lambda^d = C \frac{\lambda^{k+1}}{1 - \lambda}
\end{equation}
\end{proof}

\section{Complexity Analysis}

\begin{itemize}
    \item \textbf{Time Complexity:} $\mathcal{O}(k \cdot (d_h^2 + d_h d_x))$ per chunk update; total sequence time $\mathcal{O}(T \cdot (d_h^2 + d_h d_x))$, matching full BPTT.
    \item \textbf{Space Complexity:} $\mathcal{O}(k \cdot d_h)$ memory footprint, strictly independent of total sequence length $T$.
\end{itemize}

\section{Exactness and Error Analysis}

When the temporal dependency horizon $\Delta t$ of the underlying physical system is smaller than the truncation window $k$, TBPTT is asymptotically exact. For $\Delta t > k$, TBPTT fails to capture long-range credit assignment.

\section{Worked Numerical Example}

Let $k=1, d_x=1, d_h=1, d_y=1$. Parameters: $W_x = [[1.0]], W_h = [[0.5]], W_y = [[1.0]]$.
Sequence: $x_1 = 1.0, x_2 = 1.0$; targets $y_1^* = 0.0, y_2^* = 0.0$.

\textbf{Chunk 1 ($t=1$, $h_{\text{init}}=0.0$):}
\begin{itemize}
    \item $h_1 = \tanh(1.0) = 0.7616$, $y_1 = 0.7616$.
    \item Loss $= 0.5(0.7616)^2 = 0.2900$.
    \item Backward gives $\frac{\partial \mathcal{L}^{(1)}}{\partial W_h} = \delta_1 h_0 = 0.0$.
    \item Carry state $h_{\text{carry}} = 0.7616$.
\end{itemize}

\textbf{Chunk 2 ($t=2$, detached $h_{\text{init}}=0.7616$):}
\begin{itemize}
    \item $h_2 = \tanh(0.5 \times 0.7616 + 1.0) = \tanh(1.3808) = 0.8807$, $y_2 = 0.8807$.
    \item Loss $= 0.5(0.8807)^2 = 0.3878$.
    \item Backward gives $\delta_2 = 0.8807(1 - 0.8807^2) = 0.1976$.
    \item $\frac{\partial \mathcal{L}^{(2)}}{\partial W_h} = \delta_2 \times h_{\text{init}} = 0.1976 \times 0.7616 = 0.1505$.
    \item Backward halts; does not propagate back into Chunk 1.
\end{itemize}

\section{Numerical Considerations and Edge Cases}

\begin{itemize}
    \item \textbf{Sequence Boundaries:} At the true boundary between independent sequence episodes (e.g. document transitions), $h_{\text{carry}}$ must be explicitly zeroed: $h_{\text{carry}} \gets \mathbf{0}$.
    \item \textbf{State Detach Mechanics:} Failure to detach $h_{\text{carry}}$ in autograd frameworks creates memory leaks, eventually triggering out-of-memory crashes.
\end{itemize}

\section{References}
\begin{enumerate}
    \item Williams, R. J., & Peng, J. (1990). An efficient gradient-based algorithm for on-line training of recurrent network paths. \emph{Neural Computation}, 2(4), 490--501.
    \item Sutskever, I. (2013). \emph{Training recurrent neural networks}. PhD thesis, University of Toronto.
\end{enumerate}

\end{document}
